% R124B prospective mathematics, formulated after R124A and before B's draw.
\section{Paid values and certified policy improvement}
\label{sec:valuebridge124b}

The first experiment distinguishes two failures of the proposed
information calculation: a saturated one-anchor certificate and a
two-anchor menu that never presents a distinct action. The latter
problem requires a change in action generation. We therefore formulate
a second, separately frozen experiment. Its state draw is independent
of the first experiment, and its inference is conditional on this new
preparation.
All information still comes from paid evaluations of the original
Bellman objective; the actor, future controller, information filtration,
capacity and economic objective remain the same.

The additional baseline is a cheap conventional calculation from paid
values and original curvature. Its interpolation inequality is classical.
The methodological question is whether generating actions from these
levels changes the economic decision, and whether the more expensive
joint value--marginal certificate improves upon this baseline after all
costs are charged. A neural encoding has no exclusive access to either
calculation.

\subsection{A strong-convexity bound from two paid values}

Fix the original state and write $q=Q_x$, with curvature $\mu>0$ on
$[0,A]$. For two distinct paid actions $s<t$, retain
$W_s\le q(s)\le U_s$ and $W_t\le q(t)\le U_t$. Put $d=t-s$.
For $b=s+ud$, $u\in[0,1]$, define
\begin{equation}\label{eq:valuebridge124b}
 C_{st}(b)=(1-u)U_s+uU_t-\frac\mu2u(1-u)d^2.
\end{equation}

\begin{proposition}[A paid-value action bound]
\label{prop:valuebridge124b}
Throughout $[s,t]$, $q(b)\le C_{st}(b)$. Over the information class
containing only these two level intervals and curvature between
$\mu$ and $L\ge\mu$, this is the exact pointwise upper bound.
The minimizer of this upper bound on the bracket is
\begin{equation}\label{eq:bridgeproposal124b}
 u_{st}^*=\operatorname{clip}_{[0,1]}
       \left(\frac12-\frac{U_t-U_s}{\mu(t-s)^2}\right),
 \qquad b_{st}^*=s+(t-s)u_{st}^*.
\end{equation}
The midpoint is another explicitly feasible proposal. All formulas
are rational for rational inputs; neither a risk query nor a nonlinear
optimization is required. The upper bound is evaluated at the
\emph{realized} action. It may only be used for a pair that brackets
that action.
\end{proposition}

The sharpness assertion is for the specified two-level information,
not for all derivative intervals, raw supports and original formulas.
Those additional data can reduce its upper bound. The bridge supplies
a strong conventional comparator precisely because it uses the known
curvature and all available pairs, without a representation restriction.

\subsection{The common action menu and complete certificates}

Let $a_F$ be the completed fused service's feasible action, and retain
its entire original paid ledger. If that ledger contains at least two
distinct actions, the new fixed menu contains the rounded pairwise
minimizers \eqref{eq:bridgeproposal124b}, the pairwise midpoints, and
the six neighbors
\[
 a_F\pm A/16,\qquad a_F\pm A/64,\qquad a_F\pm A/256,
\]
clipped to the original interval. The prescribed feasible rounding
and duplicate convention are applied before certification. This menu
is generated once, using $a_F$ and the original paid ledger, and is
identical for the bridge and joint comparisons. The later selection
of a bridge action does not regenerate the menu. Retention at the
current incumbent is always available.

No interpolation upper is used outside its paid pair. In particular,
if rounding moves a proposed action outside one pair, that pair is
omitted when evaluating the actual action. The existing fused upper
remains available. When the paid endpoints lie on the same downward
rounding grid, rounding any point between them remains between them,
but correctness does not rely on this extra property.

Write $\underline q(b)$ and $U_F(b)$ for the original-q lower and
upper bounds furnished by the paid supports and fused derivative
envelopes. Let $B\le\min q$ be the common original-optimum lower
bound, and let $C_F(a,b)$ be the fused upper bound on $q(b)-q(a)$.
All these quantities use the original paid information, including
the same raw-support bounds as before. Set
\begin{align}
 U_B(b)&=\min\left\{U_F(b),\,
                  \min_{s\le b\le t} C_{st}(b)\right\},\notag\\
 C_B(a,b)&=\min\{C_F(a,b),U_B(b)-\underline q(a)\},
                      \qquad b\ne a ,\notag\\
 U_B^q(a,b)&=\min\{U_B(b),U_B(a)+C_B(a,b)\}.
 \label{eq:bridgecertificate124b}
\end{align}
An empty inner minimum is $+\infty$. For $b=a$ set the contrast
to zero by identity. The incumbent lower in this expression must
minorize $q(a)$; an arbitrary lower bound on $J(a)$ is not silently
treated as a lower bound on $q(a)$.

For the joint stage, insert its current incumbent $a_B$ and the same
candidate $b$ as unobserved nodes in the joint interpolation program.
This does not purchase a value at either point. Let $D(a_B,b)$ be
the rational dual upper contrast from
Theorem~\ref{thm:dualcertificate124}, using the unchanged paid
ledger and its valid finite boxes. Put
\begin{align}
 C_J(a_B,b)&=\min\{C_B(a_B,b),D(a_B,b)\},\notag\\
 U_J^q(a_B,b)&=\min\{U_B(b),U_B(a_B)+C_J(a_B,b)\}.
 \label{eq:jointbridgecertificate124b}
\end{align}
Thus the more expensive stage also receives the cheap value bound.
It does not compare itself with a conventional baseline deprived of
the same paid levels or feasible actions.

\begin{theorem}[Original-policy improvement through the two stages]
\label{thm:stagedpolicy124b}
Suppose the original future controller is unchanged and satisfies
$J(b)=q(b)+\eta(b)$ with $0\le\eta\le\delta$. No smoothness or
convexity of $\eta$ is required. For either stage choose a valid
original-q contrast upper $C(a,b)$ and candidate upper $U^q(b)$
from \eqref{eq:bridgecertificate124b} or
\eqref{eq:jointbridgecertificate124b}. At every distinct candidate,
\begin{align}
 J(b)-J(a)&\le C(a,b)+\delta,\notag\\
 \underline q(b)&\le J(b)\le U^q(b)+\delta,\notag\\
 0\le J(b)-\min q&\le U^q(b)+\delta-B.
 \label{eq:stagedqualification124b}
\end{align}
Accept only if the first upper bound is strictly negative, the
last upper bound is at most $\tau$, and the absolute interval width
is at most $\omega$, after all outward conversions at the realized
feasible action. Select the largest certified gain with the fixed
action/index tie rule; otherwise retain the current qualified policy.
Then every original state has
\begin{equation}\label{eq:stagedorder124b}
 V_0\le J(a_J)\le J(a_B)\le J(a_F)\le V_0+\tau,
\end{equation}
and each returned policy retains the required objective-width
qualification. The sum of the accepted stage gain lower bounds
minorizes $J(a_F)-J(a_J)$. A recoverable failure in the joint stage
retains $a_B$; it does not discard a previously certified bridge gain.
\end{theorem}

All numerical lower, upper, contrast and optimum bounds in this
theorem must concern the same original state, risk and continuation.
The future allowance is added once to a changed-action comparison
and once to its absolute policy upper, serving different predicates;
it is never subtracted from an unidentified policy upper. Retention
has exactly zero gain. A gain certificate alone does not establish
the objective-width condition.

The second experiment's restriction to at least two paid anchors is
a declared menu design. If a fused service has moved away from its
sole paid anchor, that anchor need no longer equal its incumbent;
the single-anchor saturation theorem does not automatically apply
to the new origin. This design therefore makes no claim that every
one-anchor continuation is information-optimal.

\paragraph{Information and complete work.}
The new stages add no original-risk query. They do add pair enumeration,
rounding, full candidate certification, rejected candidates, and
possibly joint optimization and rational verification. With $p$ paid
records, the unmerged extra menu has at most $2\binom p2+6$ entries;
equal-anchor pairs are omitted. Direct evaluation of the pair bounds
at $m$ candidates takes at most $m\binom p2$ scalar evaluations.
The implemented stage also recomputes the origin bound for each
candidate, giving at most $2m\binom p2$ pair evaluations per stage,
in addition to the retained fused and raw-support work. Each
standalone service independently incurs its whole acquisition and
incumbent-service bill. Mathematical dominance in
\eqref{eq:stagedorder124b} is not dominance in original objective
plus computation cost, and it is not a neural-specific approximation
theorem. Both questions require the separately frozen complete-cost
comparison with every conventional representation.
